% Part: incompleteness
% Chapter: incompleteness-provability
% Section: rosser-thm

\documentclass[../../../include/open-logic-section]{subfiles}

\begin{document}

% SEGMENT OLP-0316-S01

\olfileid{inc}{inp}{ros}

\olsection{ரோஸரின் தேற்றம்}

கோடலின் நிறுவலை மாற்றி, “$\omega$-முரணற்றது” என்பதற்குப்
பதிலாக “முரணற்றது” என்று மட்டும் வைத்து இன்னும் வலிமையான முடிவைப்
பெற முடியுமா? ரோஸர் கண்டுபிடித்த ஒரு தந்திரத்தைப் பயன்படுத்தி,
இதற்கான விடை ஆம். $\OProv[\Th{T}](y)$ க்கு பதிலாக
$\ORProv_T(y)$ என்ற “மாற்றிய” !!{derivability} பயனிலையைப்
பயன்படுத்துவதே ரோஸரின் தந்திரம்.

\begin{thm}
\ollabel{thm:rosser}
  $\Th{T}$ ஐ, $\Th{Q}$ ஐ நீட்டிக்கும் எந்த முரணற்ற, !!{axiomatizable} கோட்பாடான
$\Th{T}$ உம் முழுமையானதல்ல.
\end{thm}

\begin{proof}
$\OProv[\Th{T}](y)$ என்பது
$\lexists[x][\OPrf[\Th{T}](x, y)]$ என்று வரையறுக்கப்பட்டது என்பதை
நினைவுகூர்க. இங்கு $\OPrf[\Th{T}](x, y)$ தீர்மானிக்கத்தக்க ஓர்
உறவைப் பிரதிநிதித்துவப்படுத்துகிறது: $x$ ஆனது $y$ என்ற கோடல்
எண்ணுள்ள !!{sentence} ஒன்றின் !!a{derivation} இன் கோடல் எண்ணாக
இருந்தால், எனில் மட்டுமே அந்த உறவு பொருந்தும். $x$ மற்றும்~$y$
இடையிலான மற்றோர் உறவும் தீர்மானிக்கத்தக்கது: $x$ ஆனது $y$ என்ற
கோடல் எண்ணுள்ள வாக்கியத்தின் \emph{மறுப்புரை}யின் கோடல் எண்ணாக
இருந்தால், எனில் மட்டுமே அது பொருந்தும். $\fn{not}(x)$ என்பதைப்
பின்வருமாறு செயல்படும் முதன்மை மீள்வரையறைச் சார்பாகக் கொள்க:
$x$ ஆனது $!A$ என்ற வாய்பாட்டின் குறியீடாக இருந்தால்,
$\fn{not}(x)$ ஆனது $\lnot !A$ இன் குறியீடாகும். அப்போது
$\Refut[\Th{T}](x, y)$ என்பது
$\Prf[\Th{T}](x, \fn{not}(y))$ பொருந்தினால், எனில் மட்டுமே
பொருந்தும். இதைப் பிரதிநிதித்துவப்படுத்தும் வாய்பாட்டை
$\ORefut[\Th{T}](x, y)$ என்று குறிக்கலாம். ஆகவே,
$\Th{T} \Proves \lnot !A$ என்றும் $\delta$ அதற்குரிய
!!{derivation} என்றும் இருந்தால்,
$\Th{Q} \Proves \ORefut[\Th{T}](\gn{\delta}, \gn{!A})$.
இப்போது $\ORProv[\Th{T}](y)$ ஐப் பின்வருமாறு வரையறுப்போம்:
\[
\lexists[x][(\OPrf[\Th{T}](x,y) \land \lforall[z][(z < x \lif \lnot
  \ORefut[\Th{T}](z,y))])].
\]
தோராயமாக,
“$y$ க்கு $\Th{T}$ இல் ஒரு நிறுவல் உள்ளது; மேலும் $y$ க்கான அதைவிடச்
சிறிய மறுப்புரை எதுவும் இல்லை” என்பதே
$\ORProv[\Th{T}](y)$ இன் பொருள். $\Th{T}$ முரணற்றது என்று
கருதும்போது, $\ORProv[\Th{T}](y)$ மற்றும்
$\OProv[\Th{T}](y)$ ஆகியவை ஒரே எண்களுக்கு உண்மையாகும். ஆனால்
$\Th{T}$ இல் \emph{நிறுவத்தக்கதன்மை} என்ற நோக்கில் (மெய்மைக்கும்
நிறுவத்தக்கதன்மைக்கும் வேறுபாடு உண்டு என்பதை இப்போது அறிவோம்) இவை
  வேறு பண்புகளைக் கொண்டிருக்கின்றன. $\Th{T}$ \emph{முரணாக} இருந்தால், இரண்டும்
ஒரே எண்களுக்கு \emph{உண்மையாகாது}!{} ($\ORProv[\Th{T}](y)$ என்பது
  பொதுவாக “$y$ ரோஸர் நிறுவத்தக்கது” என்று வாசிக்கப்படுகிறது. ஆனால்
  முரணான கோட்பாடுகளில் நிறுவத்தக்க !!{sentence}s ரோஸர்
நிறுவத்தக்கவையாக இல்லாமலும் இருப்பதால், இது குழப்பமாக இருக்கலாம்.
  குழப்பத்தைத் தவிர்க்க, இதை “$y$ ஷ்மோவிக்கத்தக்கது” என்று வாசிக்கலாம்.)

நிலைப்புள்ளித் துணைத்தேற்றத்தின்படி, பின்வருமாறு ஒரு வாய்பாடு
$!R_\Th{T}$ உள்ளது:
\begin{equation}
  \Th{Q} \Proves !R_\Th{T} \liff \lnot \ORProv[\Th{T}](\gn{!R_\Th{T}}).
  \ollabel{RT}
\end{equation}
\olref[1in]{thm:first-incompleteness} இன் நிறுவலுக்கு மாறாக,
$\Th{T}$ முரணற்றது என்றால் $\Th{T}$ ஆனது $!R_\Th{T}$ ஐ !!{derive}
செய்யாது என்றும், $\Th{T}$ ஆனது $\lnot !R_\Th{T}$ ஐயும்
!!{derive} செய்யாது என்றும்
இங்கு கூறுகிறோம். (வேறு சொற்களில், $\omega$-முரணற்ற தன்மை என்ற
கருதுகோள் தேவையில்லை.)

முதலில் $\Th{T} \Proves/ !R_{\Th{T}}$ என்பதை காட்டுவோம். அது
அப்படிச் செய்யும் என்று கருதுக; அப்போது~$T$ இலிருந்து
$!R_{\Th{T}}$ க்கான !!a{derivation} ஒன்று உள்ளது; அதன் கோடல் எண்
$n$ என்க. $\OPrf[\Th{T}]$ ஆனது $\Prf[\Th{T}]$ ஐ $\Th{Q}$ இல்
பிரதிநிதித்துவப்படுத்துவதால்,
$\Th{Q} \Proves
\OPrf[\Th{T}](\num{n}, \gn{!R_{\Th{T}}})$. மேலும் ஒவ்வொரு $k<n$ க்கும்,
$k$ என்பது $\lnot !R_{\Th{T}}$ இன் !!a{derivation} இன் கோடல் எண்
அல்ல; ஏனெனில் $\Th{T}$ முரணற்றது. ஆகவே ஒவ்வொரு $k<n$ க்கும்,
$\Th{Q} \Proves \lnot
\ORefut[\Th{T}](\num{k}, \gn{!R_{\Th{T}}})$. \olref[req][min]{lem:less-nsucc}
இன்படி,
$\Th{Q} \Proves \lforall[z][(z < \num{n} \lif \lnot
\ORefut[\Th{T}](z,\gn{!R_{\Th{T}}}))]$. எனவே
\[
\Th{Q} \Proves \lexists[x][(\OPrf[\Th{T}](x,\gn{!R_{\Th{T}}}) \land \lforall[z][(z
    < x \lif \lnot \ORefut[\Th{T}](z,\gn{!R_{\Th{T}}}))])],
\]
இது $\ORProv[\Th{T}](\gn{!R_{\Th{T}}})$ என்பதன் வரையறையே. \olref{RT}
  இன்படி $\Th{Q} \Proves \lnot !R_{\Th{T}}$. $\Th{T}$ ஆனது $\Th{Q}$
  ஐ நீட்டிப்பதால், $\Th{T} \Proves \lnot !R_{\Th{T}}$.
ஆனால் $\Th{T} \Proves !R_{\Th{T}}$ என்று நாம் கருதியதால், $\Th{T}$
முரணாகிவிடும்; இது தேற்றத்தின் கருதுகோளுக்கு முரணானது.

இப்போது $\Th{T} \Proves/ \lnot !R_{\Th{T}}$ என்பதை காட்டுவோம்.
மீண்டும், அது அவ்வாறு செய்கிறது என்றும்,
$n$ என்பது $\lnot !R_{\Th{T}}$ க்கான !!a{derivation} இன் கோடல் எண்
என்றும் கருதுக. அப்போது
$\Refut[\Th{T}](n, \Gn{!R_{\Th{T}}})$ உண்மை; மேலும்
$\ORefut[\Th{T}]$ ஆனது $\Refut[\Th{T}]$ ஐ $\Th{Q}$ இல்
பிரதிநிதித்துவப்படுத்துவதால்,
$\Th{Q} \Proves \ORefut[\Th{T}](\num{n}, \gn{!R_{\Th{T}}})$.
இதனால் $\Th{T}$ முரணாகிவிடும்: அது அந்த மறுதலை ஏற்கெனவே
நிறுவுவதோடு $!R_{\Th{T}}$ ஐயும் !!{derive} செய்யும் என்பதை
மீண்டும் காட்டுவோம். இந்த நீட்சி காரணமாக
\begin{align*}
\Th{Q} & \Proves !R_{\Th{T}} \liff \lnot \ORProv[\Th{T}](\gn{!R_{\Th{T}}}),
\intertext{மேலும் $\Th{T}$ ஆனது~$\Th{Q}$ ஐ நீட்டிப்பதால், பின்வருவதை காட்டினால் போதும்:}
\Th{Q} & \Proves \lnot \ORProv[\Th{T}](\gn{!R_{\Th{T}}}).
\end{align*}
$\lnot \ORProv[\Th{T}](\gn{!R_{\Th{T}}})$ என்ற !!{sentence}, அதாவது
\begin{align*}
  \lnot & \lexists[x][(\OPrf[\Th{T}](x,\gn{!R_{\Th{T}}}) \land
    \lforall[z][(z < x \lif \lnot \ORefut[\Th{T}](z,\gn{!R_{\Th{T}}}))])],
  \intertext{என்பது தருக்க ரீதியாக}
  & \lforall[x][(\OPrf[\Th{T}](x,\gn{!R_{\Th{T}}}) \lif
    \lexists[z][(z < x \land \ORefut[\Th{T}](z,\gn{!R_{\Th{T}}}))])].
\end{align*}
இதை $\Th{Q}$ !!{derive}s செய்யும் உண்மைகளைப் பயன்படுத்தி,
தருக்கத்துடன் முறைசாரா முறையில் வாதிடுவோம். $x$ தன்னிச்சையானது என்றும்
$\OPrf[\Th{T}](x,\gn{!R_{\Th{T}}})$ என்றும் கருதுக. ஏற்கனவே
$\Th{T} \Proves/ !R_{\Th{T}}$ என்று அறிந்துள்ளோம்; எனவே ஒவ்வொரு $k$
க்கும் $\Th{Q} \Proves \lnot
\OPrf[\Th{T}](\num{k}, \gn{!R_{\Th{T}}})$. இதிலிருந்து ஒவ்வொரு $k$
க்கும் $\eq/[x][\num{k}]$ என்று பெறுகிறோம். குறிப்பாக,
(a) $\eq/[x][\num{n}]$. மேலும்
$\lnot(\eq[x][\num{0}] \lor \eq[x][\num{1}]
\lor \dots \lor \eq[x][\num{n-1}])$; எனவே
\olref[req][min]{lem:less-nsucc} இன்படி, (b)
$\lnot(x < \num{n})$. \olref[req][min]{lem:trichotomy} இன்படி
$\num{n} < x$. $\Th{Q} \Proves
\ORefut[\Th{T}](\num{n}, \gn{!R_{\Th{T}}})$ என்பதால்,
$\num{n} < x \land \ORefut[\Th{T}](\num{n}, \gn{!R_{\Th{T}}})$.
இதிலிருந்து
$\lexists[z][(z < x \land
\ORefut[\Th{T}](z,\gn{!R_{\Th{T}}}))]$ கிடைக்கிறது. $x$ தன்னிச்சையானது
என்பதால், தேவையான
\[
\lforall[x][(\OPrf[\Th{T}](x,\gn{!R_{\Th{T}}}) \lif
  \lexists[z][(z < x \land \ORefut[\Th{T}](z,\gn{!R_{\Th{T}}}))])].
\]
என்ற தேவையான வாய்பாடு கிடைக்கிறது.
\end{proof}

\begin{prob}
இயல் எண்களின் $A$, $B$ என்ற இரு கணங்கள், $A \subseteq X$ என்றும்
$B \subseteq \Complement{X}$ என்றும் அமையும் !!{decidable} கணம் $X$
எதுவும் இல்லாவிட்டால் \emph{கணிக்கத்தக்கவாறு பிரிக்கவியலாதவை} என்று
அழைக்கப்படுகின்றன ($\Complement{X}$ என்பது
$\Nat \setminus X$ என்ற நிரப்பு~$X$). $\Th{Q}$ இன் முரணற்ற
!!{axiomatizable} நீட்சியான $\Th{T}$ ஐ எடுத்துக்கொள்க. $A$ என்பது
$\Th{T}$ இல் நிறுவத்தக்க !!{sentence}s இன் கோடல் எண்களின் கணம் என்றும்,
$B$ என்பது $\Th{T}$ இல் மறுக்கத்தக்க வாக்கியங்களின் கோடல் எண்களின்
கணம் என்றும் கருதுக. $A$, $B$ கணிக்கத்தக்கவாறு பிரிக்கவியலாதவை
என்பதை நிறுவுக.
\end{prob}

\end{document}
